% YAAC Another Awesome CV LaTeX Template
%
% This template has been downloaded from:
% https://github.com/darwiin/yaac-another-awesome-cv
%
% Author:
% Christophe Roger
%
% Template license:
% CC BY-SA 4.0 (https://creativecommons.org/licenses/by-sa/4.0/)
%Section: Work Experience at the top
\sectionTitle{Professional Experience}{\faSuitcase}
%\renewcommand{\labelitemi}{$\bullet$}
\begin{experiences}
	\experience
	{Current}	{Embedded Software Engineer}{DroneShield - Australia}
	{Nov 2023}	{
		\begin{itemize}
			\item Delivered angle-of-arrival estimation and the conducted calibration behind it for a four-antenna direction-finding system (400\,MHz--7\,GHz). Designed simultaneous multi-channel injection through a characterised splitter, producing a per-channel correction table that removes amplitude and phase imbalance and makes bearing accuracy repeatable unit to unit.
			\item Cut analog front-end calibration from over 12 hours to 40 minutes -- a \textbf{20x} improvement that removed the single largest bottleneck in production. Required a full rearchitecture and coordination across production, hardware and test engineering.
			\item Built Python (Dash/Plotly) analysis tooling for frequency sweeps and calibration results. Stitched S2P measurements across 11 filter banks into a single validated front-end response.
			% TODO(pat): Sensitivity check -- specific bands, array geometry and calibration
			% architecture on a shipping counter-UAS product may be commercially sensitive.
			% Check the first two bullets against what DroneShield says publicly. Fallback:
			% "multi-antenna, wideband receive array" instead of the frequency range and
			% element count.
		\end{itemize}
	}
	{Embedded C++, Python, Rust, RF Calibration, Direction Finding, FPGA/RFSoC, Qt, gRPC, Embedded Linux, numpy/scipy, Dash/Plotly, POSIX Shell, Jira, GitLab}
	\emptySeparator
	\experience
		{May 2025}	{Co-founder \& CTO}{Borne Clothing - Australia}
		{Sep 2019}	{
			\begin{itemize}
				\item Co-founded and grew a social-enterprise clothing business to a successful exit. Donated \$30,000+ to the Against Malaria Foundation -- 5,000 families protected from mosquito-borne disease for three years.
			\end{itemize}
		}
		{Business Operations, Web Development, Lean Growth}
	\emptySeparator
	\experience
	{Sep 2023}	{Embedded Systems Engineer}{ResusRight - Australia}
	{May 2022}	{
		\begin{itemize}
			\item \textbf{Led firmware development of a manufacturing test jig}, managing a small team and setting timelines that delivered on schedule and unblocked board production.
			\item Architected a Qt application on an embedded Linux SOM, including system architecture and the BLE stack. Developed firmware for ``Nemo'', a clinical resuscitation monitor now used on real patients (NXP and Silicon Labs MCUs, BLE and serial stacks).
			\item Took over ``Juno'' firmware (ESP32, BLE) and designed its Flutter companion app with Firebase auth/database and iOS/Android Bluetooth stacks.
			\item Worked to IEC 62304, IEC 60601 and ISO 13485 -- safety-critical, traceable firmware in a startup moving at pace.
		\end{itemize}
	}
	{Embedded C/C++, Embedded Linux, Qt, NXP \mu C, Silicon Labs \mu C, Espressif \mu C, Bluetooth LE, IEC62304, IEC60601, ISO13485, Flutter, Firebase, Python, GitHub}
\emptySeparator
	\experience
		{May 2022}	{Mechatronic Engineer}{CORDEL (formerly Airsight) - Australia}
		{Apr 2018}	{
			\begin{itemize}
				% TODO(pat): VERIFY AND SHARPEN THE ArduPilot/GNSS BULLETS. Claude wrote them
				% from your description. Put the real specifics in: which airframes, what you
				% personally built vs. operated, and what the GNSS research was trying to answer.
				\item \textbf{Worked hands-on with ArduPilot and PX4 flight stacks} across the company's survey UAV fleet -- airframe bring-up, flight controller configuration and tuning, mission planning and field flight operations.
				\item Contributed to \textbf{GNSS navigation research}, characterising RTK positioning accuracy and the behaviour of GNSS-based navigation under real field conditions.
				\item \textbf{Implemented an Extended Kalman Filter fusing IMU, RTK GNSS and LiDAR} for pose estimation and localisation of a survey platform in unstructured outdoor environments. Required Bayesian filtering, linear algebra and careful handling of dropout and noise in real field data.
				\item Developed pose and velocity estimation from optical flow in a C++/OpenCV pipeline. Spent a year in New York working largely unsupervised on 3D point extraction from 2D camera footage on a moving train.
				\item Implemented a pseudo-RTOS to arbitrate commands over the serial port and designed the status indicator board -- PCB, C and AVR assembly firmware.
			\end{itemize}
		}
		{ArduPilot, PX4, UAV Flight Operations, GNSS/RTK, Mission Planning, C++, Python, OpenCV, Extended Kalman Filter, Sensor Fusion, LiDAR, RTK GNSS, Point Clouds, Embedded C, AVR Assembly, RTOS, KiCad, Rust, Jira, Bitbucket}
\end{experiences}
